% element: 10-s057:e04
\BeleskeID{10-s057}
Kada su oba ulaza niska, ulazni tranzistor je u direktnom zasićenju. $T_2$ i $T_3$ nemaju uslove da provode, dok ih gornji totem-pole stepen ima, pa izlaz ostaje visok.

Ako je samo jedan ulaz nizak, stanje se ne menja: odgovarajući emiter zadržava $T_1$ u zasićenju. Tek kada su oba ulaza visoka, $T_1$ prelazi u inverzni aktivni režim i obezbeđuje pobudu narednih tranzistora. $T_2$ i $T_3$ su tada u zasićenju i izlaz je nizak. To daje NI funkciju. Karakteristika razlikuje blaži početni pad, sa naznačenim nagibom približno $-1.6$, od strme prelazne zone.
